\subsection{COMPACT SUBSETS OF R}

THEOREM 2 (Borel-Lebesgue). For a subset of the real line R to be compact it is necessary and sufficient that it be closed and bounded.

\begin{enumerate}
\def\labelenumi{\arabic{enumi})}
\item
  The condition is necessary. Let A be a compact subset of R and let \(a\) be a real number \(>0\) . The set A is closed (Chapter I, § 9, no. 3, Proposition 4) and there exists a finite number of points \(x_{i}\) ( \(1 \leqslant i \leqslant n\) ) of R such that A is contained in the union of the neighbourhoods \([x_{i} - a, x_{i} + a]\) (Chapter I, § 9, no. 3). Let \(b\) be the maximum of the numbers \(|x_{i}|\) ; then we have \(A \subset [-b - a, b + a]\) .
\item
  The condition is sufficient. It is enough to show that every interval \([-a, +a]\) (a \textgreater{} 0) is compact, and since this interval is a closed subset of a complete uniform space, it is enough to show that, for each b \textgreater{} 0, we can cover \([-a, +a]\) by a finite number of intervals of the form \([x - b, x + b]\) (Chapter II, § 4, no. 2, Corollary to Theorem 3). Now, let n be an integer \textgreater{} 0 such that a \textless{} nb; if \(x \in [-a, +a]\) and if m is the largest integer (positive or negative) such that \(mb \leqslant x\) , then we have \(-n \leqslant m \leqslant n\) and \(mb \leqslant x \leqslant (m + 1)b\) . Hence the \(2n + 1\) intervals \([(k - 1)b, (k + 1)b]\) ( \(-n \leqslant k \leqslant n\) ) form a covering of the required type.
\end{enumerate}

COROLLARY I. A subset of the real line R is relatively compact if and only if it is bounded.

COROLLARY 2. The real line is a locally compact space and is not compact.

Remark. Theorem 2 is often referred to as the ``Heine-Borel Theorem''; see the Historical Notes to Chapters II and IV.

